% 自测题9：填空9.2.1—9.2.7。印刷204／PDF212页。
\subsection{自测题9：薄膜、劈尖与双缝填空}

\sourcetest{9.2.1}{204}{212}
\problemtext 单色光在折射率$n_1$介质中波长为$\lambda$，
垂直入射折射率$n>n_1$、厚度$e$薄膜，
膜两侧环境均为$n_1$；求上下表面反射光相遇的相位差。
\workings
真空波长$\lambda_0=n_1\lambda$。
上反射有$\pi$变相、下反射无反射半波损失；
膜内双程光程$2ne$对应相位$4\pi ne/\lambda_0$。
故
\[
\boxed{\Delta\phi=\frac{4\pi ne}{n_1\lambda}+\pi
\pmod{2\pi}}.
\]
\crosscheck 若定义两束相位差为上束减下束，
公式整体取负或把$\pi$换成$-\pi$，
但干涉可观测强度不变。

\sourcetest{9.2.2}{204}{212}
\problemtext 两相干点光源至A点等距离。
在S$_1$到A的光路插入厚度$e$、折射率$n$的玻璃片，
求A点相位差；若$\lambda=500$nm、$n=1.5$，
A恰为第四级亮纹中心，求$e$。
\workings
玻璃片替代等厚空气，单程光程增加$(n-1)e$，
\[
\boxed{|\Delta\phi|=\frac{2\pi(n-1)e}{\lambda}}.
\]
第四级亮纹要求$(n-1)e=4\lambda$，
\[
\boxed{e=\frac{4(500\,\mathrm{nm})}{0.5}
=4000\,\mathrm{nm}=4.0\,\mu\mathrm m}.
\]
\crosscheck 本装置光线\emph{单程}经过玻璃，
不能误用迈克耳孙仪往返的因子2。

\sourcetest{9.2.3}{204}{212}
\problemtext 杨氏双缝在空气中条纹间距为1.00mm，
整个装置置于折射率$4/3$的水中，其他几何尺寸固定，
求新的条纹间距。
\workings
入水后频率不变，波长变为原来的$1/n$，
条纹间距$\Delta x=D\lambda/d$也缩短$1/n$：
\[
\boxed{\Delta x_{\rm water}=1.00/(4/3)\,\mathrm{mm}
=0.750\,\mathrm{mm}}.
\]
\crosscheck 题设固定几何距离；若实际透镜同时浸水并改变焦距，
需另计光学成像条件，但原题的杨氏屏距模型未涉及透镜。

\sourcetest{9.2.4}{204}{212}
\problemtext 空气中透明劈尖的楔角$\alpha=1.0\times10^{-4}$rad，
垂直入射波长700nm单色光，邻近同类反射干涉纹距0.25cm，
求膜折射率。
\workings
近轴小楔角$t(x)\simeq\alpha x$，相邻同类纹
厚度差$\lambda/(2n)$，
\[
\boxed{n=\frac{\lambda}{2\alpha\Delta x}
=\frac{700\times10^{-9}}{2(10^{-4})(2.5\times10^{-3})}
=1.40}.
\]
\crosscheck cm与mm换算：0.25cm=2.5mm，不能误作0.25mm。

\sourcetest{9.2.5}{204}{212}
\problemtext 牛顿环装置原为空气，某一第$k$级暗环半径$r_1$；
透镜与板之间填入折射率$n$的液体后
同一级暗环半径变为$r_2$，求$n$。
\workings
同一暗环级数和镜面曲率半径下
$r_k^2=k\lambda_0R/n$，
因此
\[
\boxed{n=\left(\frac{r_1}{r_2}\right)^2}.
\]
\crosscheck 这里环级数应保持相同，
且默认填入液体不会改变反射相位的级数归属。

\sourcetest{9.2.6}{204}{212}
\problemtext 迈克耳孙仪反射镜平移0.620mm，
共有2300条纹移过参考点，求单色光波长（nm）。
\workings
镜面平移$\Delta d$使往返光程差改变$2\Delta d$，
每一$\lambda$改变对应一条纹迁移：
\[
\boxed{\lambda=\frac{2\Delta d}{N}
=\frac{2(0.620\times10^{-3})}{2300}
\simeq539.1\,\mathrm{nm}}.
\]
\crosscheck 漏掉因子2会使结果变成约270nm。

\sourcetest{9.2.7}{204}{212}
\problemtext 两束相干光各自光强都是$I_0$，
问干涉区内所能出现的最大光强。
\workings
两束光场的振幅各为$\sqrt{I_0}$（忽略相同比例常数），
同相叠加时振幅翻倍，
\[
\boxed{I_{\max}=(\sqrt{I_0}+\sqrt{I_0})^2=4I_0}.
\]
\crosscheck 不能只把两个\emph{强度}相加为$2I_0$，
因为本题是相干振幅叠加。
